\section{Limitations}
\label{sec:limitations}

\textbf{No convergence claim.} Every number reported is the accuracy of the
best-validation checkpoint within a fixed budget, not a best-achievable
accuracy. The budgets are sized to be generously past the validation peak,
and we disclose per set size whether the validation signal plateaued within
its budget (Section~\ref{sec:results_curve}); for any split where it did not,
the curve point is a lower bound on what longer training might reach.

\textbf{Selection and evaluation peek at curves.} Two protocol choices
inflate absolute numbers. Both are shared group protocol, so we keep them
and account for them here. The downstream metric is the best test top-1
across the 100 probe epochs, which selects on the test curve; the shared
evaluation script does not yet persist the final-epoch accuracy, so the gap
is unquantified in this draft and reporting it is planned for the final
version. The shipped checkpoint is
the maximum of a smoothed validation signal, and such a maximum is an
upward-biased estimate~\cite{cawley2010overfitting}; smoothing reduces the
bias but does not eliminate it, and we do not quantify what remains. Both
effects inflate every curve point in the same direction. Their magnitude
need not be identical across $N$, so they are a caveat on the curve's exact
shape as well as on absolute values.

\textbf{The validation probe is sparse.} The in-domain selection signal is a
200-way $k$NN classification with 25 gallery images per class --- far sparser
than published kNN protocols (DINO's ImageNet gallery has roughly 1\,300 per
class~\cite{caron2021dino}). Its absolute accuracy is accordingly low, and we
use it only as a relative signal for ranking checkpoints within a run; we did
not validate that this regime ranks checkpoints as a denser probe would.

\textbf{One backbone, one transfer pair, one evaluation regime.} All claims
are scoped to ViT-Tiny/8 pre-trained on Tiny-ImageNet and probed linearly on
CIFAR-10. The backbone is fixed by the shared group protocol, which controls
it as a variable but means no claim extends to other architectures or sizes.
Fine-tuning --- the regime under which the family-level data-robustness
ranking was established~\cite{elnouby2021largescale} --- is not evaluated,
and a single downstream task cannot show whether the learned invariances
help or hurt elsewhere; a multi-task probe suite is future work.

\textbf{One completed seed; no data resampling.} The numbers in this draft
come from a single pre-training seed (42); seeds 43 and 44 are in progress,
and every cross-split difference smaller than a typical seed band ---
including the 1.8-point dip itself --- should be read with that reservation
(in the earlier three-seed campaign, per-split seed ranges spanned 0.5 to
2.9 points). The splits are also never resampled: the nesting attributes
curve differences to $N$ alone, but the spread will not capture sensitivity
to \emph{which} images a small split contains. The decoupling analysis of
Section~\ref{sec:results_decoupling} rests on the kNN diagnostic, not the
linear probe; the probe-level oracle comparison is pending. Finally, the
campaign spans three minor revisions of the training code (a diagnostic
checkpoint-saving addition entered mid-campaign, uncommitted for the three
largest splits); the diff does not touch the training computation, but
strict single-commit provenance holds only per split, not across the
curve.

\textbf{Transferred schedule evidence and a-priori choices.} The
WSD-schedule findings we build on --- constant-rate competitiveness,
cooldown shape and length --- come from language-model pre-training at
scales far above ours~\cite{hagele2024scaling,hu2024minicpm}. Their transfer
to Barlow Twins on 1k--100k images is an assumption of this study. Several
recipe constants ($\lambda = 10^{-2}$, projector width 1024, crop lower
bound 0.25, disabled blur, 500 warmup steps) were set a priori and not
ablated. Finally, absolute
accuracies shifted slightly under a framework-and-GPU change during
development --- the random-initialization floor moved by about 1.2 points
--- so all reported comparisons use numbers from a single software and
hardware environment.
